% =====================================================================
%  Chapter 8: Measuring the Spatial Economy
%  Source: EconAnalysis.tex (Overleaf) — complete content
% =====================================================================

\chapter{Measuring the Spatial Economy}
\label{chap_EconAnalysis}

\section{Modelling and Measuring Economic Activity}
\label{sec_EconActivity}

Quantitative measures of economic structure serve two purposes in regional
analysis. First, they describe the industrial composition of a local economy
relative to a reference geography (usually the national average), revealing
its specialisations and dependencies. Second, they form the empirical basis
for testing theories of agglomeration, cluster formation, and regional growth.

\subsection{Measures of Spatial Concentration and Specialisation}

\subsubsection{Locational Measures}

\paragraph{Location Quotient}\index{location quotient}

The \textbf{location quotient} (LQ) is the standard measure of employment
distribution that controls for the size of the region. The relative
concentration of industry $i$ in region $j$ is:

\begin{tcolorbox}[keyresult]
\begin{equation}
  LQ_i = \frac{E_{ij}/E_{in}}{E_j/E_n}
  \label{eq_LQ}
\end{equation}
\end{tcolorbox}

\noindent where $E_{ij}$ is employment in industry $i$ in region $j$,
$E_j$ is total employment in region $j$, $E_{in}$ is national employment
in industry $i$, and $E_n$ is total national employment.

An $LQ > 1$ indicates above-average concentration of industry $i$ in region
$j$; industries with $LQ > 1.25$ are generally viewed as constituting the
region's export-oriented \textbf{economic base}\index{economic base}.
An $LQ < 1$ indicates a sector that serves primarily local demand.

\paragraph{Horizontal Clustering}\index{clustering!horizontal}

\textbf{Horizontal clustering} (HC) is an alternative measure that accounts
for possible agglomeration effects in terms of the absolute number of jobs
in a specific industry. It is the arithmetic equivalent of a bubble chart
that accounts for an industry's gravitational pull over and above what is
expected at the national level:
\begin{equation}
  HC_i = E_{ij} - \hat{E}_{ij}
\end{equation}
where $\hat{E}_{ij} = E_j \times (E_{in}/E_n)$ is the expected employment
level obtained by setting $\widehat{LQ}_i = 1$.

\subsubsection{Shift-Share Analysis}\index{shift-share analysis}

Shift-share analysis decomposes regional employment growth into three
components:
\begin{equation}
  \Delta E_{ij} = \underbrace{E_{ij}\,\frac{\Delta E_n}{E_n}}_{\text{National growth}}
  + \underbrace{E_{ij}\!\left(\frac{\Delta E_{in}}{E_{in}} - \frac{\Delta E_n}{E_n}\right)}_{\text{Industry mix}}
  + \underbrace{E_{ij}\!\left(\frac{\Delta E_{ij}}{E_{ij}} - \frac{\Delta E_{in}}{E_{in}}\right)}_{\text{Competitive effect}}
\end{equation}

The \textbf{national growth} component attributes growth to the overall
national trend. The \textbf{industry mix} (or proportionality) component
asks whether the region specialises in fast- or slow-growing industries
nationally. The \textbf{competitive} (or differential) component captures
whether local industries outperform their national counterparts---the
residual attributable to location-specific factors.

% TODO: Add worked example with BLS QCEW data

\subsubsection{Agglomeration Measures}

Location measures generally ignore geographic dispersion and the uneven
distribution of employment within the subregions of an area under study.
\textbf{Agglomeration measures}\index{agglomeration} are designed to overcome
this shortcoming.

\paragraph{Herfindahl-Hirschman Index}\index{Herfindahl-Hirschman index}

The \textbf{HHI} measures employment concentration across subregions of an
MSA:
\begin{equation}
  HHI_i = \sum_{j=1}^{n} s_j^2
\end{equation}
where $s_j$ is industry $i$'s LQ at the level of subregion $j$ relative to
the MSA, and $n$ is the number of counties within the MSA.\footnote{For the
Blacksburg MSA, $n=4$ (Giles, Montgomery, Pulaski counties and the city of
Radford).} The index equals 1 under absolute concentration and $1/n$ under
uniform dispersion.

\paragraph{Locational Gini Coefficient}\index{Gini coefficient!locational}

The \textbf{locational Gini coefficient} (LGC) measures industrial
concentration within a region more rigorously than the HHI:
\begin{equation}
  LGC_i = \frac{\sum_{i=1}^{n}\sum_{j=1}^{n}|x_i - x_j|}{2n(n-1)\mu}
\end{equation}
where $x_i$ and $x_j$ are the LQs of industries $i$ and $j$ in each of the
subregions, $\mu$ is the mean LQ in the MSA, and $n$ is the number of
subregions.

\paragraph{Ellison-Glaeser Index}\index{Ellison-Glaeser index}

The \textbf{Ellison-Glaeser} (EG) index \citep{Ellison1997} measures how
industrial concentration patterns differ from a situation where firms locate
randomly:
\begin{equation}
  EG_i = \frac{\displaystyle\sum_{j=1}^{n}(s_j - x_j)^2
               - (1-\sum_{j=1}^{n}x_j^2)\sum_{k=1}^{m}z_k^2}
              {\displaystyle(1-\sum_{j=1}^{n}x_j^2)(1-\sum_{k=1}^{m}z_k^2)}
\end{equation}
where $s_j$ is the employment share of subregion $j$ in industry $i$;
$x_j$ is the overall employment share of subregion $j$; $z_k$ is the share
of 6-digit NAICS subsector establishments in sector $k$; and $m$ is the
number of 6-digit subsectors.

The EG index accounts for both geographic concentration and industrial
concentration (plant-level returns to scale), yielding a ``dartboard''
measure: if the observed concentration exceeds what would arise if firms
chose locations randomly, the excess is attributed to agglomeration
economies.


\section{Measuring Inequality}
\label{sec_Inequality}

% TODO: Lorenz curves, Gini coefficient, Theil index
% TODO: Spatial decomposition: between-region vs. within-region inequality

Inequality measurement in a spatial context requires decomposing overall
inequality into components attributable to differences \emph{across} regions
and differences \emph{within} regions. The Theil index is particularly useful
because it admits an exact additive decomposition:
\begin{equation}
  T = T_\text{between} + T_\text{within}
\end{equation}

% TODO: Worked example with ACS income data by census tract


\section{Spatial Concentration and Agglomeration}
\label{sec_SpatialConc}

% TODO: Connect EG index to NEG theory; add empirical evidence on
%       US industrial clusters (Porter 2003, Glaeser et al. 1992)


\section{Economic Impact Analysis}
\label{sec_EconImpact}

\subsection{Input-Output Modelling}
\label{subsec_IO}

The \textbf{Leontief input-output model}\index{input-output model} describes
the interdependencies among sectors of an economy. Let $A$ be the
\textbf{technical coefficients matrix}\index{technical coefficients matrix},
where $a_{ij}$ is the amount of good $i$ required to produce one unit of
good $j$. The model is:
\begin{equation}
  \mathbf{x} = A\mathbf{x} + \mathbf{f}
  \qquad\Longrightarrow\qquad
  (I - A)\mathbf{x} = \mathbf{f}
\end{equation}
where $\mathbf{x}$ is the vector of total output by sector and $\mathbf{f}$
is the vector of final demand. The solution:

\begin{tcolorbox}[keyresult]
\begin{equation}
  \mathbf{x} = (I - A)^{-1}\mathbf{f} = L\mathbf{f}
\end{equation}
where $L = (I-A)^{-1}$ is the \textbf{Leontief inverse}\index{Leontief
inverse} (total requirements matrix).
\end{tcolorbox}

Each element $l_{ij}$ of $L$ gives the total output of sector $i$ required,
directly and indirectly, to deliver one unit of sector $j$'s output to
final demand.

\paragraph{Regional multipliers.}

The \textbf{output multiplier} for sector $j$ is $\sum_i l_{ij}$ (column
sum of $L$): the total output generated across all sectors per unit of
final demand for sector $j$. BEA's RIMS II
multipliers\index{RIMS II multipliers} are derived from the national I-O
accounts supplemented with regional employment data.

\noindent\textbf{Open-source I-O analysis.} BEA's publicly available
detailed use tables, combined with the \texttt{leontief} package in R,
permit full regional multiplier estimation without proprietary software:

\begin{lstlisting}[language=R, caption={Open I-O analysis with BEA data}]
library(leontief)
# Read BEA Use Table (downloaded from bea.gov)
use_table <- read.csv("bea_use_table.csv", row.names = 1)
A <- technical_coefficients(use_table)  # compute A matrix
L <- leontief_inverse(A)                # compute L = (I-A)^{-1}
multipliers <- colSums(L)               # output multipliers
\end{lstlisting}

\subsection{Applied Tools}
\label{subsec_AppliedIO}

\subsubsection{RIMS II}\index{RIMS II}

The Bureau of Economic Analysis (BEA) Regional Input-Output Modelling System
(RIMS II) provides multipliers for any area (state, MSA, county) for any
industry defined in the BEA national I-O accounts. Multipliers are available
for output, earnings, and employment. They are derived from two primary data
sources: the national I-O accounts and the BEA Regional Economic Accounts.

\subsubsection{IMPLAN}\index{IMPLAN}

IMPLAN is a proprietary social accounting matrix (SAM) model that extends
the standard I-O framework by adding household income and government accounts.
It is widely used in planning practice for economic impact assessments.
The key advantage over RIMS II is its spatial flexibility (any custom
geographic boundary) and its ability to model household spending patterns.
The key disadvantage is cost and the ``black box'' nature of proprietary
multipliers.

\subsubsection{FedFIT}\index{FedFIT}

FedFIT (Federal Footprint in the Territories) is a BEA tool for estimating
the economic contribution of federal government spending in US territories
and states.

\paragraph{A note on interpreting multipliers.}

Common mistakes in economic impact analysis include: (i)
\textbf{double-counting}---each sector's contribution should be computed once;
(ii) \textbf{ignoring displacement}---new activity may displace existing
activity rather than generating net new output; (iii) \textbf{treating I-O
estimates as causal}---I-O multipliers describe input requirements under a
fixed technology, not causal effects of policy interventions. For causal
analysis of economic development interventions, quasi-experimental
identification strategies (Chapter~\ref{chap_BeyondOLS}) are required.

\renewcommand\bibname{Further reading}
\begin{subappendices}
\bibliographystyle{econometrica}
\bibliography{Literature/OverLord}
\IfFileExists{Literature/EconAnalysis.bbl}{\chapter{Economic Analysis}\label{chap_EconAnalysis}

\section{Modelling and measuring economic activity}
\subsection{Measures of spatial concentration and specialisation}
The various quantitative measures of economic concentration\index{measures of economic concentration} are defined in detail in this section.

\subsubsection{Locational measures}

\paragraph{Location quotient}
The location quotient\index{location quotient|see{measures of economic concentration}}\index{measures of economic concentration!location quotient} (LQ) is the standard measure of employment distribution that controls for the size of the region. The relative
concentration of industry $i$ in region $j$ is defined as

\begin{eqnarray}
LQ_{i} = \frac{(E_{ij}/E_{in})}{(E_{j}/E_{n})}
\end{eqnarray}

\noindent where, $E_{ij}$ is employment in industry $i$ in region $j$, $E_{j}$ is total employment in region $j$, $E_{in}$ is national employment in industry $i$, and $E_{n}$ is total national employment. Thus, a LQ of greater than one indicates that there is an above average proportion of employment in a given industry in a given region. Industries with an LQ above 1.25 are generally viewed
as constituting the export-oriented economic base.

\paragraph{Horizontal clustering}
Horizontal clustering\index{clustering!horizontal|see{measures of economic concentration}}\index{measures of economic concentration!horizontal clustering} (HC) is an alternative measure that accounts for possible agglomeration effects in terms of the numbers of jobs in a specific industry. As such, it is the arithmetic equivalent to a bubble chart\index{bubble chart} which accounts for an industry's gravitational pull over and above what is expected to occur at the national level. HC can then be written as

\begin{eqnarray}
HC_{i} = E_{ij} - \widehat{E}_{ij},
\end{eqnarray}

\noindent where $\widehat{E}_{ij} = E_{j}\times \frac{E_{in}}{E_{n}}$ by setting $\widehat{LQ}_{i}=1$.

\subsubsection{Shift-share analysis}


\subsubsection{Agglomeration measures}
Location measures generally ignore geographic dispersion and uneven distribution of employment within subregions of an area under study.
Agglomeration measures\index{agglomeration}\index{agglomeration|seealso{measures of economic concentration}} are designed to overcome this shortcoming.

\paragraph{Herfindahl-Hirschman Index}
In the context of industry employment concentration within an MSA the simplest such measure is the Herfindahl-Hirschman index\index{Herfindahl-Hirschman index|see{measures of economic concentration}}\index{measures of economic concentration!Herfindahl-Hirschman index} (HHI). It is defined as

\begin{eqnarray}
HHI_{i} = \sum^{n}_{i=1}s_{i}^{2},
\end{eqnarray}

\noindent where $s_{i}$ is simply industry $i$'s LQ at the level of the subregion compared to the MSA and $n$ is the number of counties
within the MSA.\footnote{In the context of MSA-level data, counties are commonly assumed to be the subregions. For the Blacksburg MSA,
we would set $n=4$ since it includes Giles, Montgomery, Pulaski counties and the city of Radford.} The index is equal to 1 if there is absolute concentration and it takes a value of $\frac{1}{n}$ if employment in the industry is equally dispersed across the MSA.

\paragraph{Locational Gini Coefficient}
The locational Gini coefficient\index{Gini coefficient!locational} (LGC) also accounts for agglomeration and concentration within a specific region, but in a most sophisticated way than the HHI. It is defined as

\begin{eqnarray}
LGC_{i} =
\frac{\sum_{i=1}^{n}\sum_{j=1}^{n}|x_{i}-x_{j}|}{2n(n-1)\mu},
\end{eqnarray}

\noindent where $x_{i}$ and $x_{j}$ are the LQs of industries $i$ and $j$ in each of the subregions, $\mu$ is the mean of the LQs in the MSA and $n$ is the number of counties.

\paragraph{Ellison-Glaeser Index}
The Ellison-Glaeser index\index{Ellison-Glaeser index|see{measures of economic concentration}}\index{measures of economic concentration!Ellison-Glaeser index} (EG) measures how industrial concentration patterns differ from a situation where firms in a purely random manner. The index is defined as follows

\begin{equation}
EG_{i} = \frac{\sum_{i=1}^{n}(s_{i}-x_{i})^2 - (1-
\sum^{n}_{i=1}x^{2}_{i})\sum^{m}_{j=1}z^{2}_{j}}{(1-
\sum^{n}_{i=1}x^{2}_{i})(1- \sum^{m}_{j=1}z^{2}_{j})},
\end{equation}

\noindent where $s_{i}$, $x_{i}$ and $n$ have the same definitions as above; $z_{i}$ is the share of the 6-digit NAICS subsector
establishments in sector $j$ and $m$ is the number of 6-digit subsectors.


\section{Measuring inequality}

\section{Spatial concentration and agglomeration}

\section{Economic impact analysis}
\subsection{Input-output modelling}
\subsection{Applied tools}
\subsubsection{RIMS II}
\subsubsection{IMPLAN}
\subsubsection{FedFIT}

\renewcommand\bibname{Further reading}
\begin{subappendices}
\bibliographystyle{econometrica}
\bibliography{Literature/OverLord}
\IfFileExists{Literature/EconAnalysis.bbl}{\chapter{Economic Analysis}\label{chap_EconAnalysis}

\section{Modelling and measuring economic activity}
\subsection{Measures of spatial concentration and specialisation}
The various quantitative measures of economic concentration\index{measures of economic concentration} are defined in detail in this section.

\subsubsection{Locational measures}

\paragraph{Location quotient}
The location quotient\index{location quotient|see{measures of economic concentration}}\index{measures of economic concentration!location quotient} (LQ) is the standard measure of employment distribution that controls for the size of the region. The relative
concentration of industry $i$ in region $j$ is defined as

\begin{eqnarray}
LQ_{i} = \frac{(E_{ij}/E_{in})}{(E_{j}/E_{n})}
\end{eqnarray}

\noindent where, $E_{ij}$ is employment in industry $i$ in region $j$, $E_{j}$ is total employment in region $j$, $E_{in}$ is national employment in industry $i$, and $E_{n}$ is total national employment. Thus, a LQ of greater than one indicates that there is an above average proportion of employment in a given industry in a given region. Industries with an LQ above 1.25 are generally viewed
as constituting the export-oriented economic base.

\paragraph{Horizontal clustering}
Horizontal clustering\index{clustering!horizontal|see{measures of economic concentration}}\index{measures of economic concentration!horizontal clustering} (HC) is an alternative measure that accounts for possible agglomeration effects in terms of the numbers of jobs in a specific industry. As such, it is the arithmetic equivalent to a bubble chart\index{bubble chart} which accounts for an industry's gravitational pull over and above what is expected to occur at the national level. HC can then be written as

\begin{eqnarray}
HC_{i} = E_{ij} - \widehat{E}_{ij},
\end{eqnarray}

\noindent where $\widehat{E}_{ij} = E_{j}\times \frac{E_{in}}{E_{n}}$ by setting $\widehat{LQ}_{i}=1$.

\subsubsection{Shift-share analysis}


\subsubsection{Agglomeration measures}
Location measures generally ignore geographic dispersion and uneven distribution of employment within subregions of an area under study.
Agglomeration measures\index{agglomeration}\index{agglomeration|seealso{measures of economic concentration}} are designed to overcome this shortcoming.

\paragraph{Herfindahl-Hirschman Index}
In the context of industry employment concentration within an MSA the simplest such measure is the Herfindahl-Hirschman index\index{Herfindahl-Hirschman index|see{measures of economic concentration}}\index{measures of economic concentration!Herfindahl-Hirschman index} (HHI). It is defined as

\begin{eqnarray}
HHI_{i} = \sum^{n}_{i=1}s_{i}^{2},
\end{eqnarray}

\noindent where $s_{i}$ is simply industry $i$'s LQ at the level of the subregion compared to the MSA and $n$ is the number of counties
within the MSA.\footnote{In the context of MSA-level data, counties are commonly assumed to be the subregions. For the Blacksburg MSA,
we would set $n=4$ since it includes Giles, Montgomery, Pulaski counties and the city of Radford.} The index is equal to 1 if there is absolute concentration and it takes a value of $\frac{1}{n}$ if employment in the industry is equally dispersed across the MSA.

\paragraph{Locational Gini Coefficient}
The locational Gini coefficient\index{Gini coefficient!locational} (LGC) also accounts for agglomeration and concentration within a specific region, but in a most sophisticated way than the HHI. It is defined as

\begin{eqnarray}
LGC_{i} =
\frac{\sum_{i=1}^{n}\sum_{j=1}^{n}|x_{i}-x_{j}|}{2n(n-1)\mu},
\end{eqnarray}

\noindent where $x_{i}$ and $x_{j}$ are the LQs of industries $i$ and $j$ in each of the subregions, $\mu$ is the mean of the LQs in the MSA and $n$ is the number of counties.

\paragraph{Ellison-Glaeser Index}
The Ellison-Glaeser index\index{Ellison-Glaeser index|see{measures of economic concentration}}\index{measures of economic concentration!Ellison-Glaeser index} (EG) measures how industrial concentration patterns differ from a situation where firms in a purely random manner. The index is defined as follows

\begin{equation}
EG_{i} = \frac{\sum_{i=1}^{n}(s_{i}-x_{i})^2 - (1-
\sum^{n}_{i=1}x^{2}_{i})\sum^{m}_{j=1}z^{2}_{j}}{(1-
\sum^{n}_{i=1}x^{2}_{i})(1- \sum^{m}_{j=1}z^{2}_{j})},
\end{equation}

\noindent where $s_{i}$, $x_{i}$ and $n$ have the same definitions as above; $z_{i}$ is the share of the 6-digit NAICS subsector
establishments in sector $j$ and $m$ is the number of 6-digit subsectors.


\section{Measuring inequality}

\section{Spatial concentration and agglomeration}

\section{Economic impact analysis}
\subsection{Input-output modelling}
\subsection{Applied tools}
\subsubsection{RIMS II}
\subsubsection{IMPLAN}
\subsubsection{FedFIT}

\renewcommand\bibname{Further reading}
\begin{subappendices}
\bibliographystyle{econometrica}
\bibliography{Literature/OverLord}
\IfFileExists{Literature/EconAnalysis.bbl}{\chapter{Economic Analysis}\label{chap_EconAnalysis}

\section{Modelling and measuring economic activity}
\subsection{Measures of spatial concentration and specialisation}
The various quantitative measures of economic concentration\index{measures of economic concentration} are defined in detail in this section.

\subsubsection{Locational measures}

\paragraph{Location quotient}
The location quotient\index{location quotient|see{measures of economic concentration}}\index{measures of economic concentration!location quotient} (LQ) is the standard measure of employment distribution that controls for the size of the region. The relative
concentration of industry $i$ in region $j$ is defined as

\begin{eqnarray}
LQ_{i} = \frac{(E_{ij}/E_{in})}{(E_{j}/E_{n})}
\end{eqnarray}

\noindent where, $E_{ij}$ is employment in industry $i$ in region $j$, $E_{j}$ is total employment in region $j$, $E_{in}$ is national employment in industry $i$, and $E_{n}$ is total national employment. Thus, a LQ of greater than one indicates that there is an above average proportion of employment in a given industry in a given region. Industries with an LQ above 1.25 are generally viewed
as constituting the export-oriented economic base.

\paragraph{Horizontal clustering}
Horizontal clustering\index{clustering!horizontal|see{measures of economic concentration}}\index{measures of economic concentration!horizontal clustering} (HC) is an alternative measure that accounts for possible agglomeration effects in terms of the numbers of jobs in a specific industry. As such, it is the arithmetic equivalent to a bubble chart\index{bubble chart} which accounts for an industry's gravitational pull over and above what is expected to occur at the national level. HC can then be written as

\begin{eqnarray}
HC_{i} = E_{ij} - \widehat{E}_{ij},
\end{eqnarray}

\noindent where $\widehat{E}_{ij} = E_{j}\times \frac{E_{in}}{E_{n}}$ by setting $\widehat{LQ}_{i}=1$.

\subsubsection{Shift-share analysis}


\subsubsection{Agglomeration measures}
Location measures generally ignore geographic dispersion and uneven distribution of employment within subregions of an area under study.
Agglomeration measures\index{agglomeration}\index{agglomeration|seealso{measures of economic concentration}} are designed to overcome this shortcoming.

\paragraph{Herfindahl-Hirschman Index}
In the context of industry employment concentration within an MSA the simplest such measure is the Herfindahl-Hirschman index\index{Herfindahl-Hirschman index|see{measures of economic concentration}}\index{measures of economic concentration!Herfindahl-Hirschman index} (HHI). It is defined as

\begin{eqnarray}
HHI_{i} = \sum^{n}_{i=1}s_{i}^{2},
\end{eqnarray}

\noindent where $s_{i}$ is simply industry $i$'s LQ at the level of the subregion compared to the MSA and $n$ is the number of counties
within the MSA.\footnote{In the context of MSA-level data, counties are commonly assumed to be the subregions. For the Blacksburg MSA,
we would set $n=4$ since it includes Giles, Montgomery, Pulaski counties and the city of Radford.} The index is equal to 1 if there is absolute concentration and it takes a value of $\frac{1}{n}$ if employment in the industry is equally dispersed across the MSA.

\paragraph{Locational Gini Coefficient}
The locational Gini coefficient\index{Gini coefficient!locational} (LGC) also accounts for agglomeration and concentration within a specific region, but in a most sophisticated way than the HHI. It is defined as

\begin{eqnarray}
LGC_{i} =
\frac{\sum_{i=1}^{n}\sum_{j=1}^{n}|x_{i}-x_{j}|}{2n(n-1)\mu},
\end{eqnarray}

\noindent where $x_{i}$ and $x_{j}$ are the LQs of industries $i$ and $j$ in each of the subregions, $\mu$ is the mean of the LQs in the MSA and $n$ is the number of counties.

\paragraph{Ellison-Glaeser Index}
The Ellison-Glaeser index\index{Ellison-Glaeser index|see{measures of economic concentration}}\index{measures of economic concentration!Ellison-Glaeser index} (EG) measures how industrial concentration patterns differ from a situation where firms in a purely random manner. The index is defined as follows

\begin{equation}
EG_{i} = \frac{\sum_{i=1}^{n}(s_{i}-x_{i})^2 - (1-
\sum^{n}_{i=1}x^{2}_{i})\sum^{m}_{j=1}z^{2}_{j}}{(1-
\sum^{n}_{i=1}x^{2}_{i})(1- \sum^{m}_{j=1}z^{2}_{j})},
\end{equation}

\noindent where $s_{i}$, $x_{i}$ and $n$ have the same definitions as above; $z_{i}$ is the share of the 6-digit NAICS subsector
establishments in sector $j$ and $m$ is the number of 6-digit subsectors.


\section{Measuring inequality}

\section{Spatial concentration and agglomeration}

\section{Economic impact analysis}
\subsection{Input-output modelling}
\subsection{Applied tools}
\subsubsection{RIMS II}
\subsubsection{IMPLAN}
\subsubsection{FedFIT}

\renewcommand\bibname{Further reading}
\begin{subappendices}
\bibliographystyle{econometrica}
\bibliography{Literature/OverLord}
\IfFileExists{Literature/EconAnalysis.bbl}{\input{Literature/EconAnalysis.bbl}}{\typeout{Need to run bibtex}}
\end{subappendices} }{\typeout{Need to run bibtex}}
\end{subappendices} }{\typeout{Need to run bibtex}}
\end{subappendices} }{\typeout{Need bibtex on EconAnalysis}}
\end{subappendices}
